L'étiquette d'un médicament fournit l'information
<<\textbf{Ibuprofène.... 400 mg}>>.

On dissout un comprimé contenant l'ibuprofène selon un protocole bien défini
afin d'obtenir une solution aqueuse $\mathrm{(S)}$ d'ibuprofène de volume
$V_{S}=\qty{100}{\mL}$.

Pour vérifier, la masse d'ibuprofène contenu dans ce comprimé, on procède à un
titrage acido-basique du volume $V_{S}$ par une solution aqueuse d'hydroxyde de
sodium $\ce{Na^{+}_{(aq)} + HO^{-}_{(aq)}}$ de concentration molaire
$C_{B}=\qty{1.94e-1}{\mol\per\L}$, en utilisant le dispositif expérimental de
la Fig.~\ref{fig:2bac_svt_national_2018_normale_chim_1}.

La Fig.~\ref{fig:2bac_svt_national_2018_normale_chim_2} donne les courbes
$\mathrm{pH}=f(V_{B})$ et
$\frac{\mathrm{d}\mathrm{pH}}{\mathrm{d}V_{B}}=g(V_{B})$ obtenues lors de ce
dosage.

\begin{minipage}{0.49\linewidth}
    \begin{figure}[H]
        \centering
        \includegraphics[width=0.75\textwidth]{2026-03-12_015233-}
        \caption{}
        \label{fig:2bac_svt_national_2018_normale_chim_1}
    \end{figure}
\end{minipage}
\hfill
\begin{minipage}{0.49\linewidth}
    \begin{figure}[H]
        \centering
        \includegraphics[width=0.8\textwidth]{2026-03-12_015243-}
        \caption{}
        \label{fig:2bac_svt_national_2018_normale_chim_2}
    \end{figure}
\end{minipage}

\begin{enumerate}
    \item Nommer les éléments du dispositif expérimental numérotés 1, 2, 3 et 4
          sur la Fig.~\ref{fig:2bac_svt_national_2018_normale_chim_1}.
    \item Parmi les courbes~(1) et~(2) de la
          Fig.~\ref{fig:2bac_svt_national_2018_normale_chim_2}, quelle est celle
          qui représente $\mathrm{pH}=f(V_{B})$~?
    \item Déterminer graphiquement la valeur du volume $V_{B,E}$ versé à
          l'équivalence.
    \item Écrire l'équation de la réaction qui a eu lieu lors du dosage sachant
          qu'elle est totale.
    \item Calculer la valeur de la quantité de matière $n_{A}$ d'ibuprofène dans
          la solution $\mathrm{(S)}$.
    \item Déduire la valeur de la masse $m$ d'ibuprofène dans le comprimé et la
          comparer à celle indiquée sur l'étiquette du médicament.
\end{enumerate}

